%%%%%%%%%%%%%%%%%%%%%%%%%%%%%%%%%%%%%%%%%%%%%%%%%%%%%%%%%%%%%%%%%%%%%%
% Lớp 11 | Dãy số. Cấp số cộng. Cấp số nhân | Cấp số nhân
% Dạng 1D2?3-5: Điều kiện để dãy số là cấp số nhân
% Số câu: 38  (sinh tự động bởi tach_cau.py; từ nay thầy sửa trực tiếp file này)
%%%%%%%%%%%%%%%%%%%%%%%%%%%%%%%%%%%%%%%%%%%%%%%%%%%%%%%%%%%%%%%%%%%%%%

%%%%%%%%%%%%%%%% TRẮC NGHIỆM 4 PHƯƠNG ÁN %%%%%%%%%%%%%%%%

% nguồn: Ok_Toan11_TracNghiem_Lan2
\begin{ex}%[1D2H3-5]%[ma:f4ecce9f]
	Xác  định $x$ dương để $2x-3$; $x$; $2x+3$ lập thành cấp số nhân.
	\choice 
		 {$x=3$}
		{\True  $x=\sqrt 3 $}
		{$x=\pm\sqrt  3 $}
		{Không có giá trị  nào của $x$}
	 \loigiai{
		$2x-3$; $x$; $2x+3$ lập thành  cấp số nhân\\
		$\Leftrightarrow{x^2}=\left(2x-3\right)\left(2x+3\right) \Leftrightarrow{x^2}=4x^2-9  \Leftrightarrow{x^2}=3 \Leftrightarrow x=\pm\sqrt 3$.\\
	Vì $x$ dương nên $x=\sqrt  3 $.}
\end{ex}

% nguồn: Ok_Toan11_TracNghiem_Lan2
\begin{ex}%[1D2H3-5]%[ma:44bf1789]
	Có bao nhiêu giá trị nguyên dương của $x$ để ba số $1$; $x$; $x+2$ theo thứ tự đó lập thành một cấp  số nhân?
	 \choice
		{$2$} 
		{\True  $1$}
		 {$3$}
		{$0$} 
	 \loigiai{
		Để $1$;  $x$; $x+2$ theo thứ tự đó lập thành một cấp số nhân thì
		 $x^2=x+2 \Leftrightarrow \heva{ &x=-1\\ &x=2.}$\\
	Vậy có  đúng $1$ số nguyên dương $x=2$.}
\end{ex}

% nguồn: Ok_Toan11_TracNghiem_Lan2
\begin{ex}%[1D2H3-5]%[ma:7a2f78ef]
	Cho cấp số nhân $(u_n)$ có số  hạng đầu $u_1=2$ và công bội $q=5$. Giá trị của $\sqrt{u_6\cdot u_8}$ bằng
	 \choice
		{\True $2\cdot  5^6$}
		{$2\cdot 5^7$} 
		{$2\cdot 5^8$} 
		{$2\cdot 5^5$} 
	 \loigiai{
	$(u_n)$ là cấp số nhân $\Rightarrow u_6u_8={u_7}^2\Rightarrow\sqrt{u_6u_8}=|u_7|=|u_1q^6|=2\cdot 5^6$. 
	 }
\end{ex}

% nguồn: Ok_Toan11_TracNghiem_Lan2
\begin{ex}%[1D2H3-5]%[ma:85a18b1f]
	Có bao nhiêu giá trị nguyên của $x$ để $3$ số $1$; $x$; $x+2$ theo  thứ tự lập thành một cấp số nhân?
	\choice 
		{$2$} 
		{\True  $1$}
		 {$3$}
		 {$0$}
	 \loigiai{
		 Để $3$ số $1$; $x$; $x+2$ theo thứ tự lập thành một cấp số nhân thì
		\[x^2=\left(x+2\right)\cdot  1\Leftrightarrow \hoac{&x=-1\\&x=2.}\]
		Vậy có đúng $1$  số nguyên dương $x=2$.
	 }
\end{ex}

% nguồn: Ok_Toan11_TracNghiem_Lan2
\begin{ex}%[1D2H3-5]%[ma:b5f42b57]
	Biết ba số $x^2$, $8$, $x$ theo thứ  tự lập thành cấp số nhân. Giá trị của $x$ là
	\choice 
		 {\True $x=4$}
		 {$x=5$}
		{$x=2$} 
		{$x=1$} 
	 \loigiai{
	Vì $x^2$, $8$, $x$ theo thứ tự lập thành  cấp số nhân.\\
	Nên $x^2\cdot x=8^2\Rightarrow x^3=64\Rightarrow x=4$. 
	} 
\end{ex}

% nguồn: Ok_Toan11_TracNghiem_Lan2
\begin{ex}%[1D2H3-5]%[ma:737bfc75]
	Tập hợp các giá trị $x$ thỏa mãn $x$, $2x$, $x+3$ theo thứ tự lập thành  một cấp số nhân là
	\choice 
		 {$\left\{0;1\right\}$}
		{$\emptyset$} 
		{\True $\left\{1\right\}$} 
		 {$\left\{0\right\}$}
	\loigiai{Gọi $q$ là công bội của  cấp số nhân.Ta có
		\[\heva{&2x=x\cdot q\\&x+3=2x\cdot q}\Leftrightarrow  \heva{&2x=x\cdot q\\&x+3=2\cdot 2x}\Rightarrow \heva{&q=2\\&x=1.}\]
		Tập hợp các giá trị $x$ thỏa mãn $x$, $2x$, $x+3$ theo thứ tự lập thành một cấp  số nhân là $x=1$.
	 }
\end{ex}

% nguồn: Ok_Toan11_TracNghiem_Lan2
\begin{ex}%[1D2H3-5]%[ma:e835615c]
	Tìm tất cả giá trị của $x$ để ba số $2x - 1$; $x$; $2x + 1$ theo thứ  tự đó lập thành một cấp số nhân.
	\choice 
		{\True  $x = \pm \dfrac{1}{\sqrt{3}}$}
		{$x  = \pm \dfrac{1}{3}$}
		 {$x = \pm \sqrt{3}$}
		{$x = \pm 3$} 
	\loigiai{ 
		Ta  có
		\[  u_1\cdot u_3=(u_1)^2\cdot q^2=(u_2)^2.\]
		Vậy ba số lập thành cấp  số nhân khi
		\begin{eqnarray*}
			&&x^2 = (2x - 1)(2x +  1)\\
			&\Leftrightarrow& x^2 =  4x^2 - 1\\
			&\Leftrightarrow& 3x^2  = 1\\
			&\Leftrightarrow& x^2 =  \dfrac{1}{3}\\
			&\Leftrightarrow & x=\pm \dfrac{\sqrt{3}}{3}. 
		\end{eqnarray*}
		Vậy  $x = \pm \dfrac{1}{\sqrt{3}}$.
	 }
\end{ex}

% nguồn: Ok_Toan11_TracNghiem_Lan2
\begin{ex}%[1D2H3-5]%[ma:293a7312]
	Tập hợp các giá trị $x$ thỏa mãn $x$,  $2x$, $x+3$ theo thứ tự lập thành một cấp số nhân là
	\choice 
		 {$\{{0;1}\}$}
		{$\varnothing $} 
		{\True $\{  1\}$}
		 {$\{ 0\}$}
	\loigiai{ 
		Gọi $q$ là công bội của cấp  số nhân.\\
		Ta  có\\
		$\heva{  &2x=x\cdot q\\
		&x+3=2x \cdot q} 
		\Leftrightarrow \heva{&2x=x\cdot  q\\
		 &x+3=2 \cdot 2x}
		 \Rightarrow \heva{&q=2\\
		&x=1.}$\\ 
	Tập hợp các giá  trị $x$ thỏa mãn $x$, $2x$, $x+3$ theo thứ tự lập thành một cấp số nhân là $\{ 1\}$.}
\end{ex}

% nguồn: Ok_Toan11_TracNghiem_Lan2
\begin{ex}%[1D2H3-5]%[ma:fadebdfb]
	Xác định $x$  để $3$ số $2x-1; x; 2x+1$ theo thứ tự lập thành một cấp số nhân:
	\choice 
		 {$x=\pm \dfrac{1}{3}$}
		 {$x=\pm \sqrt{3}$}
		 {\True $x=\pm \dfrac{1}{\sqrt{3}}$}
		{Không có giá trị  nào của $x$}
	\loigiai{ 
		Ba số: $2x-1; x; 2x+1$ theo thứ tự lập thành cấp số  nhân khi
		\begin{eqnarray*}
			(2x-1)(2x+1)=x^2 
			&\Leftrightarrow& 4x^2-1=x^2 \\ 
			 &\Leftrightarrow& 3x^2=1
			\Leftrightarrow x=\pm  \dfrac{1}{\sqrt{3}}.
		\end{eqnarray*}
	 }
\end{ex}

% nguồn: Ok_Toan11_TracNghiem_Lan2
\begin{ex}%[1D2H3-5]%[ma:2a9d4cf8]
	 Xác định $x$ để $3$ số $x - 2$; $x + 1$; $3 - x$ theo thứ tự lập thành một cấp số nhân.
	 \choice
		{\True Không có  giá trị nào của $x$}
		{$x  = \pm 1$}
		{$x  = 2$}
		 {$x = -3$}
	\loigiai{ 
		Ba số $x - 2$; $x + 1$; $3 - x$ theo thứ tự lập thành cấp  số nhân, suy ra
		\begin{eqnarray*}
			 && (x - 2)(3 - x) = (x + 1)^2\\
			 &\Leftrightarrow & 2x^2 - 3x + 7 = 0\quad \text{(phương trình vô nghiệm).}
		\end{eqnarray*}
		Vậy không  tồn tại $x$ thỏa mãn.
	 }
\end{ex}

% nguồn: Ok_Toan11_TracNghiem_Lan2
\begin{ex}%[1D2H3-5]%[ma:84a01529]
	Xác định $x$ để $3$ số $x-2;  x+1; 3-x$ theo thứ tự lập thành một cấp số nhân:
	\choice 
		{\True Không có giá trị  nào của $x$}
		{$x=\pm  1$}
		 {$x=2$}
		 {$x=-3$}
	\loigiai{ 
		Ba số $x-2; x+1; 3-x$ theo thứ tự lập thành  một cấp số nhân khi
		$$(x-2)(3-x)=(x+1)^2  \Leftrightarrow 2x^2-3x+7=0.$$
	 }
\end{ex}

% nguồn: Ok_Toan11_TracNghiem_Lan2
\begin{ex}%[1D2H3-5]%[ma:dd15b40e]
	Tìm tất cả các giá  trị của $x$ để ba số $2x-1$, $x$, $2x+1$ theo thứ tự đó lập thành một cấp số nhân.
	 \choice
		{$x=\pm\dfrac{1}{3}$} 
		 {\True $x=\pm\dfrac{1}{\sqrt 3}$}
		 {$x=\pm\sqrt 3 $}
		{$x=\pm  3$}
	\loigiai{ 
		Để ba số đó lập thành  một cấp số nhân thì\\
	$x^2=\left(2x-1\right)\left(2x+1\right)  \Leftrightarrow{x^2}=4x^2-1 \Leftrightarrow{x^2}=\dfrac{1}{3} \Leftrightarrow x=\pm\dfrac{1}{\sqrt 3}$.}
\end{ex}

% nguồn: Ok_TracNghiemLop11_Lan1
\begin{ex}%[1D2H3-5]%[ma:e1b551ac]
	Tìm tất cả giá trị của $x$ để ba số $2x$, $7x+6$, $6 x^2-24$ theo thứ tự đó lập  thành một cấp số nhân.
	\choice 
		 {$x=4$}
		{$x=5$} 
		 {\True $x=6$}
		{$x=7$} 
	\loigiai{ 
		Áp dụng  tính chất của cấp số nhân, ta có
		$$ 
		\begin{aligned}
			& (2 x)\left(6  x^2-24\right)=(7 x+6)^2 \\
			& \Leftrightarrow 12 x^3-48  x=49 x^2+84 x+36 \\
			& \Leftrightarrow 12 x^3-49 x^2-132  x-36=0
		\end{aligned}
		$$ 
		 Vì lấy $x$ nguyên nên nhận $x=6$.
	 }
\end{ex}

% nguồn: Ok_TracNghiemLop11_Lan1
\begin{ex}%[1D2H3-5]%[ma:ad663f8e]
	Tìm $b>0$ để các số $\dfrac{1}{\sqrt{2}};\sqrt{b};\sqrt{2}$ theo thứ tự đó lập thành một cấp số nhân. 
	 \choice
		 {$b=2$}
		 {$b=-1$}
		 {$b=-2$}
		 {\True $b=1$}
	 \loigiai{
	Cấp số  nhân $\dfrac{1}{\sqrt{2}};\sqrt{b};\sqrt{2}\Rightarrow \dfrac{1}{\sqrt{2}}\cdot \sqrt{2}={{(\sqrt{b})}^2}\Leftrightarrow b=1$.}
\end{ex}

% nguồn: Ok_TracNghiemLop11_Lan1
\begin{ex}%[1D2H3-5]%[ma:122474a7]
	Tìm tất cả các  giá trị của $x$ để ba số $2x-1$, $x$, $2x+1$ theo thứ tự lập thành một cấp số nhân.
	\choice 
		{\True $x=\pm  \dfrac{\sqrt{3}}{3}$}
		{$x=\pm 3$} 
		{$x=\pm  \sqrt{3}$}
		{$x=\pm  \dfrac{1}{3}$}
	\loigiai{ Ba số $2x-1$, $x$, $2x+1$ theo thứ tự lập thành một cấp số nhân nên 
		\[x^2=(2x-1)(2x+1)\Leftrightarrow 3x^2-1=0\Leftrightarrow x =\pm \dfrac{\sqrt{3}}{3}. \] 
	 }
\end{ex}

% nguồn: Ok_TracNghiemLop11_Lan1
\begin{ex}%[1D2H3-5]%[ma:eb024d74]
	Xác định $x$ để bộ ba số $2x-1$, $x$, $2x+1$ theo thứ tự  lập thành một cấp số nhân.
	 \choice
		{ \True $x=\pm \dfrac{1}{\sqrt{3}}$} 
		{$x=\pm  \sqrt{3}$}
		{Không có giá trị nào của  $x$}
		{$x=\pm \dfrac{1}{3}$} 
	\loigiai{ 
		Bộ ba số $2x-1$, $x$, $2x+1$ theo thứ tự lập thành một cấp số nhân nên  ta có $$(2x-1)(2x+1)=x^2 \Leftrightarrow 4x^2-1=x^2\Leftrightarrow x=\pm \dfrac{1}{\sqrt{3}}.$$
	} 
\end{ex}

% nguồn: Ok_TracNghiemLop11_Lan1
\begin{ex}%[1D2H3-5]%[ma:d51006a9]
	Ba số $-\sqrt{3}\,;\,x\,;-3\sqrt{3}$ theo thứ tự là ba số hạng liên tiếp của một cấp số nhân. Tìm  công bội $q$ của cấp số nhân đó.
	 \choice
		{$q=\pm  3$}
		 {$q=-\sqrt{3}$}
		{$q=3$} 
		 {\True \True $q=\pm \sqrt{3}$}
	 \loigiai{
		Do $-\sqrt{3}\,;x\,;\,-3\sqrt{3}$ là một cấp số nhân $\Rightarrow x^2=9\Leftrightarrow x=\pm  3$.\\
		Vậy  công bội của cấp số nhân là $q=\dfrac{x}{-\sqrt{3}}=\pm \sqrt{3}$.
	} 
\end{ex}

% nguồn: Ok_TracNghiemLop11_Lan1
\begin{ex}%[1D2H3-5]%[ma:4b5193c0]
	Tìm tất cả giá trị của $x$ để ba số $2x-1;\,x;\,2x+1$ theo thứ tự đó lập thành  một cấp số nhân.
	\choice 
		 {$ x=\pm \dfrac{1}{3}$}
		{\True $ x=\pm \dfrac{1}{\sqrt {3}}$} 
		{$  x=\pm \sqrt {3}$}
		{$  x=\pm 3$}
	 \loigiai{
		Vì $3$ số $2x-1;\,x;\,2x+1$  lập thành cấp số nhân nên
		$$(2x-1)(2x+1)=x^2\Leftrightarrow  3x^2=1\Leftrightarrow x^2=\dfrac{1}{3}\Leftrightarrow x=\pm \dfrac{1}{\sqrt {3}}.$$
	} 
\end{ex}

% nguồn: Ok_TracNghiemLop11_Lan1
\begin{ex}%[1D2H3-5]%[ma:ebeed853]
	Xác định $ x$ để $3$ số $2x-3$;  $x$; $2x+3$ theo thứ tự lập thành một cấp số nhân.
	\choice 
		 {$x=\pm\dfrac{1}{3}$}
		{\True $x=\pm\sqrt3$} 
		 {$x=\pm\dfrac{1}{\sqrt 3}$}
		{$x=\pm 3$} 
	 \loigiai{
		Theo tính chất cấp số nhân $\left(2x-3\right)\left(2x+3\right)=x^2\Leftrightarrow  4x^2-9=x^2\Leftrightarrow x=\pm\sqrt 3 $.
	 }
\end{ex}

% nguồn: Ok_TracNghiemLop11_Lan1
\begin{ex}%[1D2H3-5]%[ma:70cc2f96]
	Cho hàm số $f(x) = \dfrac{1}{5}x^5 + \dfrac{4}{3}x^3 - 5x  + 3$. Mệnh đề nào sau đây \textbf{sai}?
	\choice 
		 {Phương trình $f(x) = 0$ có nghiệm trên khoảng $(-1;1)$}
		{Phương trình $f(x) = 0$ có  nghiệm trên khoảng $(0;+ \infty)$}
		{\True  Hàm số đã cho gián đoạn tại $x_0 = \dfrac{1}{5}$}
		{Hàm số $f(x)$ liên tục  trên $\mathbb{R}$}
	\loigiai{ 
		 Tập xác định: $\mathscr{D} = \mathbb{R}$.
		 Ta có $f\left(\dfrac{1}{2} \right) \cdot f\left(\dfrac{4}{5} \right) < 0 \Rightarrow \exists x_0 \in \left(0;\dfrac{9}{10}\right): f(x_0) =0$. \\
		Do đó: \lq\lq Phương trình $f(x)  = 0$ có nghiệm trên khoảng $(-1;1)$\rq\rq\, và \lq\lq Phương trình $f(x) = 0$ có nghiệm trên khoảng $(0;+ \infty)$ \rq\rq\, là hai mệnh đề đúng.\\
		Hàm số $f(x)$ là hàm đa  thức nên liên tục trên tập xác định $\mathscr{D} = \mathbb{R}$.\\
		Nên mệnh đề \lq\lq Hàm số đã cho gián đoạn tại $x_0 = \dfrac{1}{5}$\rq\rq \, là mệnh đề sai. 
	 }
\end{ex}

% nguồn: Ok_TracNghiemLop11_Lan1
\begin{ex}%[1D2H3-5]%[ma:b9b44041]
	Cho dãy số $\dfrac{-1}{\sqrt{2}};\, \sqrt{b};\, \sqrt{2}$. Chọn $b$  để dãy số đã cho lập thành cấp số nhân?
	\choice 
		{  $b=-1$}
		{$b=1$} 
		 {$b=2$}
		{\True Không có giá trị nào của $b$} 
	\loigiai{ 
		Dãy số đã cho lập thành cấp số nhân  khi $\heva{&b \geq 0 \\& b=-\dfrac{1}{\sqrt{2}} \cdot \sqrt{2}=-1}$.\\
		Vậy  không có giá trị nào của $b$.
	} 
\end{ex}

% nguồn: Ok_TracNghiemLop11_Lan1
\begin{ex}%[1D2H3-5]%[ma:7abde11d]
	Tìm tất cả giá trị của $x$ để ba số $2x-1$, $x$, $2x+1$ theo  thứ tự đó lập thành một cấp số nhân.
	 \choice
		 {\True $x=\pm \dfrac{1}{\sqrt{3}}$}
		{$x=\pm  \dfrac{1}{3}$}
		 {$x=\pm \sqrt{3}$}
		{$x=\pm  3$}
	\loigiai{ 
		Ba số $2x-1$, $x$, $2x+1$ theo thứ  tự đó lập thành một cấp số nhân
		\begin{eqnarray*}
			&\Rightarrow&x^2=(2x-1)(2x+1)\\ 
			 &\Leftrightarrow&x^2=4x^2-1\\
			&\Leftrightarrow&3x^2=1\\ 
			 &\Leftrightarrow&x=\pm \dfrac{1}{\sqrt{3}}.
		\end{eqnarray*}
	} 
\end{ex}

% nguồn: Ok_TracNghiemLop11_Lan1
\begin{ex}%[1D2H3-5]%[ma:cea6410c]
	Ba số $1$, $3$, $x$ theo thứ tự lập  thành cấp số nhân. Tìm $x$.
	 \choice
		{$1$} 
		 {$3$}
		 {$5$}
		{\True  $9$}
	 \loigiai{
		Ba số $1$, $2$, $x$ theo thứ tự lập thành một cấp số nhân nên suy ra $3^2=1\cdot  x\Leftrightarrow x=9$.
	 }
\end{ex}

% nguồn: Ok_Toan11_TracNghiem_Lan2
\begin{ex}%[1D2V3-5]%[ma:acc80850]
	Giả sử $\dfrac{\sin\alpha}{6}$, $\cos\alpha $, $\tan\alpha $ theo thứ tự đó là một  cấp số nhân. Tính $\cos 2\alpha $.
	 \choice
		{$\dfrac{\sqrt  3}{2}$}
		{$-\dfrac{\sqrt 3}{2}$} 
		 {$\dfrac{1}{2}$}
		{\True $-\dfrac{1}{2}$} 
	 \loigiai{
		Điều kiện $\cos\alpha\ne  0 \Leftrightarrow\alpha\ne\dfrac{\pi}{2}+k\pi $.\\
		Theo tính chất của  cấp số nhân, ta có\\
		$\cos  ^2\alpha=\dfrac{\sin\alpha}{6} \cdot \tan\alpha \Leftrightarrow 6\cos ^2\alpha=\dfrac{\sin^2\alpha}{\cos\alpha}\\
		\Leftrightarrow 6\cos  ^3\alpha-\sin ^2\alpha=0 \Leftrightarrow 6\cos ^3\alpha+\cos ^2\alpha-1=0 \Leftrightarrow\cos\alpha=\dfrac{1}{2}$.\\
	Ta có $\cos  2\alpha=2\cos ^2\alpha-1=2 \cdot \left(\dfrac{1}{2}\right)^2-1=-\dfrac{1}{2}$.}
\end{ex}

% nguồn: Ok_Toan11_TracNghiem_Lan2
\begin{ex}%[1D2V3-5]%[ma:b79732dc]
	Biết rằng tồn tại đúng hai giá trị của tham số $m$ để phương trình $x^3-7x^2+2\left(m^2+6m\right)x-8=0$ có ba nghiệm phân biệt lập thành một  cấp số nhân. Tính tổng lập phương của hai giá trị đó.
	 \choice
		{\True  $-342$}
		 {$-216$}
		 {$344$}
		{$216$} 
	 \loigiai{
		Giả sử phương trình đã cho có $3$ nghiệm là $x_1$, $  x_2$, $ x_3$.\\
		Theo định lý Viet, tích $3$ nghiệm  $x_1 x_2 x_3=-\dfrac{d}{a}=8$.\\
		Vì ba nghiệm này lập thành một  cấp số nhân nên $x_2^2=x_1 x_3$.\\
		 Do đó ta có $x_2^3=8 \Leftrightarrow x_2=2$.\\
		Thay $x=2$ vào  phương trình ta được $4\left(m^2+6m\right)=28 \Leftrightarrow \hoac{&{m=1} \\
		&{m=-7}. 
		}$\\ 
		Theo giả thiết hai giá trị này  của $m$ đều nhận.\\
		Tổng lập phương của hai giá trị  $m$ là $1^3+\left(-7\right)^3=-342$.
	 }
\end{ex}

% nguồn: Ok_Toan11_TracNghiem_Lan2
\begin{ex}%[1D2V3-5]%[ma:78134fd7]
	Với  $x$ là số nguyên dương, ba số $2x$, $3x+3$, $5x+5$ theo thứ tự là ba số hạng liên tiếp của một cấp số nhân. Số hạng tiếp theo của cấp số nhân đó là
	 \choice
		{$-\dfrac{250}{3}$} 
		{\True $\dfrac{250}{3}$} 
		 {$250$}
		{$-250$} 
	\loigiai{ 
	Vì ba số $2x$, $3x+3$, $5x+5$ theo thứ tự là ba số  hạng liên tiếp của một cấp số nhân.\\
	Nên $2x(5x+5)=(3x+3)^2\Leftrightarrow x^2-8x-9=0\Leftrightarrow x=-1$ (loại vì $x$ là số  nguyên dương) hoặc $x=9$.\\
	Với $x=9$, suy ra cấp số  nhân có công bội $q=\dfrac{3\cdot 9+3}{2\cdot 9}=\dfrac{5}{3}$.\\
	Số hạng tiếp theo của  cấp số nhân đó là $(5\cdot 9+5)\cdot\dfrac{5}{3}=\dfrac{250}{3}$.
	} 
\end{ex}

% nguồn: Ok_Toan11_TracNghiem_Lan2
\begin{ex}%[1D2V3-5]%[ma:66e65329]
	Tìm tất cả các giá trị của tham số $m$ để phương trình sau có ba nghiệm phân biệt lập thành một cấp số  nhân: $x^3-7x^2+2\left(m^2+6m\right)x-8=0$.
	\choice 
		 {$m=-7$}
		 {$m=1$}
		{$m=-1$  hoặc $m=7$}
		{\True $m=1$ hoặc  $m=-7$}
	 \loigiai{
		\begin{itemize}
			\item Điều kiện cần: Giả sử phương trình đã cho có ba nghiệm phân biệt $x_1, x_2, x_3$ lập thành  một cấp số nhân.\\
			Theo  định lý Vi-ét, ta có $x_1 x_2 x_3=8$.\\
			Theo tính chất của cấp số nhân, ta có $x_1 x_3=x_2^2$. Suy ra ta có $x_2^3=8\Leftrightarrow  x_2=2$.\\
			Với nghiệm $x=2$, ta có $m^2+6m-7=0\Leftrightarrow \hoac{&{m=1}  \\ &{m=-7}.}$
			\item Điều kiện đủ: Với $m=1$ hoặc $m=-7$  thì $m^2+6m=7$ nên ta có phương trình \[x^3-7x^2+14x-8=0.\]
			Giải phương trình này, ta được các nghiệm là $1$, $2$, $4$. Hiển nhiên ba nghiệm này lập thành một cấp số nhân với  công bội $q=2$.
		\end{itemize}
		Vậy $m=1$ và $m=-7$ là các giá trị cần  tìm.
	} 
\end{ex}

% nguồn: Ok_Toan11_TracNghiem_Lan2
\begin{ex}%[1D2V3-5]%[ma:fbfa4e72]
	Cho năm số $a$, $b$, $c$, $d$, $e$ tạo thành một cấp số nhân theo thứ tự đó và các số đều khác $0$, biết $\dfrac{1}{a}+\dfrac{1}{b}+\dfrac{1}{c}+\dfrac{1}{d}+\dfrac{1}{e}=10$ và tổng của chúng bằng $40$. Tính giá trị $\left|S\right|$  với $S=abcde$.
	 \choice
		{$\left|S\right|=42$} 
		 {$\left|S\right|=62$}
		 {\True $\left|S\right|=32$}
		{$\left|S\right|=52$} 
	\loigiai{ 
		Gọi $q$ $\left(q\ne 0\right)$ là công bội của cấp số nhân $a$, $b$,  $c$, $d$, $e$.\\
		Khi đó $\dfrac{1}{a}$, $\dfrac{1}{b}$, $\dfrac{1}{c}$, $\dfrac{1}{d}$, $\dfrac{1}{e}$ là  cấp số nhân có công bội $\dfrac{1}{q}$.\\
		Theo đề bài  ta có
		\[\heva{&{a+b+c+d+e=40} \\ 
		&{\dfrac{1}{a}+\dfrac{1}{b}+\dfrac{1}{c}+\dfrac{1}{d}+\dfrac{1}{e}=10} 
		} \Leftrightarrow  \heva{&{a\cdot\dfrac{1-q^5}{1-q}=40} \\
		 &{\dfrac{1}{a}\cdot\dfrac{1-\left(\dfrac{1}{q} \right)^5}{1-\dfrac{1}{q}}=10}
		}  \Leftrightarrow \heva{&{a\cdot\dfrac{1-q^5}{1-q}=40} \\
		&{\dfrac{1}{a}\cdot\dfrac{q^5-1}{q^4\left(q-1\right)}=10} 
		 } \Leftrightarrow a^2q^4=4.\]
		Ta có $S=abcde=a\cdot  aq\cdot aq^2\cdot aq^3\cdot aq^4=a^5q^{10}$.\\
		Nên  $S^2=\left(a^5q^{10} \right)^2=(a^2q^4)^5=4^5$.\\
		Suy  ra $\left|S\right|=\sqrt{4^5}=32$.
	 }
\end{ex}

% nguồn: Ok_Toan11_TracNghiem_Lan2
\begin{ex}%[1D2V3-5]%[ma:076439c1]
	Có bao nhiêu giá trị thực của  tham số $m$ để phương trình $\left(x-1\right)\left(x-3\right)\left(x-m\right)=0$ có $3$ nghiệm phân biệt lập thành cấp số nhân tăng?
	\choice 
		{$4$} 
		 {\True $3$}
		{$2$} 
		 {$1$}
	\loigiai{ 
		Ta có $\left(x-1\right)\left(x-3\right)\left(x-m\right)=0\Leftrightarrow \hoac{&x=1\\&x=3 \\&x=m.}$\\ 
		Để phương trình có $3$ nghiệm phân biệt  thì $m\notin \{1; 3\}$.
		\begin{itemize}
			\item Trường hợp  1: $m < 1< 3$.\\
			 Để 3 số $m$; $1$; $3$ lập thành cấp số nhân tăng thì $m\cdot 3=1^2\Leftrightarrow m=\dfrac{1}{3}$.\\
			Cấp số nhân tăng đó  là $\dfrac{1}{3}$; $1$; $3$.
			\item Trường hợp 2: $1<  m < 3$.\\
			 Để 3 số $1$; $m$; $3$ lập thành cấp số nhân tăng thì $1\cdot 3=m^2\Leftrightarrow \hoac{&{m=\sqrt{3}} \\
			&{m=-\sqrt{3}. } 
			}$  \\
			Đối chiếu điều kiện $1< m < 3$ ta  chọn $m=\sqrt{3}$.
			Cấp số nhân tăng đó là $1$; $\sqrt{3}$; $3$. 
			\item Trường hợp 3: $1< 3< m$.\\ 
			Để $3$ số  $1$; $3$; $m$ lập thành cấp số nhân tăng thì $1\cdot m=3^2\Leftrightarrow m=9$.\\
			Cấp số nhân tăng đó là  $1$; $3$; $9$.
		\end{itemize}
		Vậy $m\in \left\{\dfrac{1}{3}; \sqrt{3};  9\right\}$ thì phương trình $\left(x-1\right)\left(x-3\right)\left(x-m\right)=0$ có $3$ nghiệm phân biệt lập thành cấp số nhân tăng.
	} 
\end{ex}


%%%%%%%%%%%%%%%% ĐÚNG/SAI %%%%%%%%%%%%%%%%

% nguồn: Ok_dungsai_lop11_Lan1
\begin{ex}%[1D2H3-5]%[ma:49d836bb]
 Cho  cấp số nhân $(u_n)$ có $u_1=\dfrac{1}{2},q=-2$.
  \choiceTF
  {\True Số hạng thứ 3 là $u_3=2$} 
   {Tổng của 5 số hạng đầu là của cấp số nhân là $\dfrac{85}{4}$}
  {Cấp số  nhân đã cho là một dãy số tăng}
  {\True Số $-1024$ là số hạng thứ 12  của cấp số nhân}
 \loigiai{ 
  \begin{itemchoice}
   \itemch Đúng. Vì  $u_3=u_1\cdot q^2=\dfrac{1}{2}\cdot (-2)^2=2$.\\
   \itemch Sai.  Vì tổng của $5$ số hạng đầu là $S_5=u_1\dfrac{1-q^5}{1-q}=\dfrac{1}{2}\cdot \dfrac{1-(-2)^5}{1-(-2)}=\dfrac{11}{2}$.\\
   \itemch Sai. Vì $u_1=\dfrac{1}{2}$, $u_2=-1$, $u_3=2$.\\ 
   \itemch Đúng. Ta có $u_n=u_1\cdot q^{n-1}=-1024$.\\ 
   Nên $u_1\cdot q^{n-1}\Leftrightarrow  \dfrac{1}{2}\cdot(-2)^{n-1}=-1024\Leftrightarrow (-2)^{n-1}=(-2)^{11}\Leftrightarrow n=12$.
 \end{itemchoice}}
\end{ex}

% nguồn: Ok_dungsai_lop11_Lan1
\begin{ex}%[1D2H3-5]%[ma:42fb6d37]
 Cho tứ giác ${A B C D}$ có bốn góc tạo thành một cấp số nhân có công bội bằng $2$. Khi đó: 
  \choiceTF
  {\True  Số đo góc nhỏ nhất bằng $24^\circ $}
  {Số đo góc lớn nhất bằng $196^\circ $} 
  {Tổng số đo góc lớn nhất với góc nhỏ nhất  bằng $220^\circ $}
  {\True Số đo góc lớn nhất trừ cho  số đo góc nhỏ nhất bằng $168^\circ $}
  \loigiai{
  Đặt ${u_1; u_2; u_3; u_4}$ theo thứ tự là số  đo bốn góc của tứ giác ${A B C D}$, gọi ${q}$ là công bội của cấp số nhân ${u_1; u_2; u_3; u_4 \Rightarrow q=2}$.\\
  Ta có:  \allowdisplaybreaks
  \begin{eqnarray*}
   \heva{&u_1+u_2+u_3+u_4=360^{\circ} \\ &q=2} &\Leftrightarrow&\heva{&u_1 \cdot \dfrac{1-q^4}{1-q}=360^{\circ}  \\ &q=2}\\
   &\Leftrightarrow&\heva{&{ u _ { 1 } \cdot  \frac { 1 - 2 ^ { 4 } }{ 1 - 2 } = 3 6 0 ^ { \circ } } \\&q = 2}\\
   & \Leftrightarrow& \heva{&u_1=24^{\circ}  \\&q=2.}
  \end{eqnarray*}
  Vậy số đo bốn góc của tứ giác ${A B C D}$  là: $24^{\circ}; 48^{\circ}; 96^{\circ}; 192^{\circ}$.
  }
\end{ex}

% nguồn: Ok_dungsai_lop11_Lan2
\begin{ex}%[1D2H3-5]%[ma:d9efeae7]
	Cho các dãy số sau đây: $u_n=(\sqrt{5})^{2n-3}$; $v_n=\dfrac{2}{n}$; $w_n=\dfrac{3^{n+1}}{2^n}$ và dãy số hữu hạn gồm  các số hạng: $16;4;1;\dfrac{1}{4};\dfrac{1}{16};\dfrac{1}{64}$. Khi đó
	\choiceTF 
		{\True $\left(u_n \right)$ là một cấp số nhân công bội  $q=5$}
		{\True $\left(v_n \right)$ không  phải là một cấp số nhân}
		{\True $\left(w_n \right)$ là một cấp số nhân có số  hạng đầu $w=\dfrac{9}{2}$}
		{Dãy số hữu hạn đã cho theo thứ  tự lập thành cấp số nhân có công bội bằng $\dfrac{1}{8}$}
	 \loigiai{
		\begin{itemchoice}
			 \itemch \textbf{Đúng.} \\
			Ta có $u_{n+1}=(\sqrt{5})^{2(n+1)-3}=(\sqrt{5})^{2n-1} \Rightarrow \dfrac{u_{n+1}}{u_n}=\dfrac{(\sqrt{5})^{2n-1}}{(\sqrt{5})^{2n-3}}=(\sqrt{5})^2=5,\forall n\in  \mathbb{N}^{*}$.\\
			Do  đó $\left(u_n \right)$ là một cấp số nhân có số hạng đầu $u_1=\dfrac{1}{\sqrt{5}}$ với công bội $q=5$.
			\itemch \textbf{Đúng.}  \\
			 Ta có $v_{n+1}=\dfrac{2}{n+1} \Rightarrow \dfrac{v_{n+1}}{v_n}=\dfrac{n}{n+1}$ (tỉ số này còn phụ thuộc vào $n$).\\
			Do đó $\left(v_n \right)$ không phải là một cấp số nhân. 
			\itemch  \textbf{Đúng.} \\
			Ta có $w_{n+1}=\dfrac{3^{n+2}}{2^{n+1}}=\dfrac{3}{2} \cdot  \dfrac{3^{n+1}}{2^n} \Rightarrow \dfrac{w_{n+1}}{w_n}=\dfrac{\dfrac{3}{2} \cdot \dfrac{3^{n+1}}{2^n}}{\dfrac{3^{n+1}}{2^n}}=\dfrac{3}{2}, \forall n\in \mathbb{N}^{*}$.\\
			Do đó $\left(w_n \right)$ là một cấp  số nhân có số hạng đầu $w_1=\dfrac{3^2}{2^1}=\dfrac{9}{2}$ với công bội $q=\dfrac{3}{2}$.
			\itemch \textbf{Đúng.} \\ 
			Đặt $u_1=16$; $u_2=4$; $u_3=1$;  $u_4=\dfrac{1}{4}$; $u_5=\dfrac{1}{16}$; $u_6=\dfrac{1}{64}$.\\
			Ta có $u_2=u_1 \cdot \dfrac{1}{4}$;  $u_3=u_2 \cdot \dfrac{1}{4}$; $u_4=u_3 \cdot \dfrac{1}{4}$; $u_5=u_4 \cdot \dfrac{1}{4}$; $u_6=u_5 \cdot \dfrac{1}{4}$.\\
			 Vì vậy dãy số hữu hạn đã cho theo thứ tự lập thành cấp số nhân có công bội bằng $\dfrac{1}{4}$.
		\end{itemchoice}
	} 
\end{ex}

% nguồn: Ok_dungsai_lop11_Lan1
\begin{ex}%[1D2B3-5]%[ma:dcfbcfe6]
 Hãy cho biết tính đúng sai  của mỗi mệnh đề dưới đây?
  \choiceTF
  {\True Với $x=18$ thì các số $2$; $6$;  $x$; $54$ theo thứ tự đó lập thành một cấp số nhân}
  {Cho các số $x+2$; $x+14$; $x+50$  theo thứ tự đó lập thành cấp số nhân. Khi đó $x$ bằng $3$}
  {Có hai giá trị nguyên dương của $x$ để ba số $1$, $x$, $x+2$ theo  thứ tự đó lập thành một cấp số nhân}
  {Gọi $S$ là tập các giá trị của $x$ để ba số $x$; $6$; $x-1$ theo thứ tự lập  thành cấp số nhân. Khi đó $S=0$}
 \loigiai{ 
  \begin{itemchoice}
   \itemch Để các số $2$; $6$; $x$; $54$ theo thứ tự đó lập  thành một cấp số nhân thì $\heva{&2\cdot x=36\\&6\cdot 54=x^2} \Leftrightarrow x=18$.
   \itemch Ta có $(x+2)(x+50)=(x+14)^2\Leftrightarrow x^2+52x+100=x^2+28x+196\Leftrightarrow  24x=96\Leftrightarrow x=4$.
   \itemch Ta  có
   \begin{eqnarray*}
    \text{Yêu cầu bài toán} \Leftrightarrow 1\cdot (x+2)=x^2 \Leftrightarrow x^2-x-2=0 \Leftrightarrow \hoac{&x=-1  &\text{ (loại)}&\\ &x=2 &\text{ (nhận)}&.}
   \end{eqnarray*}
   \itemch Ba số $x$; $6$; $x-1$ theo thứ tự lập  thành cấp số nhân
   \[\Leftrightarrow  x(x-1) = 36 \Leftrightarrow x^2 - x - 36 = 0 \Leftrightarrow \hoac{&x=\dfrac{1-\sqrt{145}}{2}\\&x=\dfrac{1+\sqrt{145}}{2}.}\]
    Do đó, $S = \left\{\dfrac{1-\sqrt{145}}{2}, \dfrac{1+\sqrt{145}}{2} \right\}$.\\
    Tổng các phần tử của $S$ bằng $1$.
  \end{itemchoice}
 } 
\end{ex}

% nguồn: Ok_dungsai_lop11_Lan1
\begin{ex}%[1D2B3-5]%[ma:153a78d8]
 Hãy cho biết tính đúng sai của mỗi mệnh đề  dưới đây?
  \choiceTF
  {Cho một cấp số nhân gồm 5 số hạng theo thứ tự là $ 4 $; $  12 $; $ 36 $; $ x $; $ 324 $. Khi đó $ x=54 $}
  {\True Cho  cấp số nhân: $\dfrac{-1}{5}$; $a$; $\dfrac{-1}{\text{125}}$. Giá trị của $a$ là $a=\pm \dfrac{1}{25}$}
  {\True Với giá trị nào của $x$ bằng $3$, ba số  $x+1$, $x+9$, $x+33$ theo thứ tự lập thành một cấp số nhân}
  {Biết bốn góc của một tứ giác lồi tạo thành một cấp số nhân và góc lớn nhất bằng $27$ lần góc nhỏ nhất.  Tính tổng số đo của góc lớn nhất và góc nhỏ nhất bằng $254^\circ$}
 \loigiai{ 
  \begin{itemchoice}
   \itemch Công bội của cấp số nhân trên là $ q=\dfrac{12}{4}=3 $. Suy ra $ x=36\cdot  q=36\cdot 3=108 $.
   \itemch Ta  có $a^2=\left(-\dfrac{1}{5}\right)\cdot\left(-\dfrac{1}{125}\right)=\dfrac{1}{625}\Leftrightarrow a=\pm \dfrac{1}{25}$.
   \itemch Ba số $x+1$, $x+9$, $x+33$ theo thứ tự  lập thành một cấp số nhân nên \[(x+1)(x+33)=(x+9)^2\Leftrightarrow x^2+34x+33=x^2+18x+81\Leftrightarrow 16x=48\Leftrightarrow x=3.\]
    Khi đó cấp số nhân là $4$, $12$, $36$.
   \itemch Gọi bốn góc của  tứ giác lồi (từ bé đến lớn) lần lượt là $x,qx,q^2x$ và $q^3x$ (với $x>0,q>1$).\\
   Theo giả  thiết, ta có $q^3x=27x \Rightarrow q^3=27 \Rightarrow q=3$.\\
    Ta có $x+3x+9x+27x=360^\circ \Rightarrow x=9^\circ$.\\
   Suy ra  tổng số đo của góc lớn nhất và góc nhỏ nhất bằng $9^\circ+27\cdot 9^\circ=252^\circ$.
  \end{itemchoice}
 } 
\end{ex}


%%%%%%%%%%%%%%%% TRẢ LỜI NGẮN %%%%%%%%%%%%%%%%

% nguồn: Ok_TraLoiNgan_Lop11_Lan2
\begin{ex}%[1D2H3-5]%[ma:f2b16878]
	Tìm $x$ để $3$ số $1+x$; $9+x$; $33+x$ theo thứ tự đó lập thành  một cấp số nhân.
	\shortans{3} 
	\loigiai{ 
		Vì $3$ số $1+x$; $9+x$; $33+x$ theo thứ tự đó lập thành  một cấp số nhân, nên ta có
		\begin{eqnarray*}
			 && \left(1+x\right)\left(33+x\right)=\left(9+x\right)^2\\
			&\Leftrightarrow& x^2+34x+33=x^2+18x+81\\ 
			 &\Leftrightarrow& 16x=48\\
			 &\Leftrightarrow& x=3.
		\end{eqnarray*}
	 }
\end{ex}

% nguồn: Ok_TraLoiNgan_Lop11_Lan2
\begin{ex}%[1D2H3-5]%[ma:93936e9d]
	Tìm $x$ để các số $2$; $8$;  $x$; $128$ theo thứ tự đó lập thành một cấp số nhân.
	 \shortans{32}
	 \loigiai{
		Cấp số nhân $2$; $8$; $x$; $128$ theo thứ tự đó sẽ là $u_1$; $u_2$;  $u_3$; $u_4$, ta có
		\[ 
		\heva{&\dfrac{u_2}{u_1}=\dfrac{u_3}{u_2}\\&\dfrac{u_3}{u_2}=\dfrac{u_4}{u_3}} \Leftrightarrow \heva{&\dfrac{8}{2}=\dfrac{x}{8}\\&\dfrac{128}{x}=\dfrac{x}{8}}\Leftrightarrow \heva{&x=32\\&x^2=1024}\Leftrightarrow \heva{&x=32\\&\hoac{&x=32\\&x=-32}} \Leftrightarrow  x=32.
		\] 
	} 
\end{ex}

% nguồn: Ok_TraLoiNgan_Lop11_Lan1
\begin{ex}%[1D2V3-5]%[ma:c6cd96d5]
	Tìm tổng tất cả các  giá trị của tham số $m$ để phương trình sau có ba nghiệm phân biệt lập thành một cấp số nhân $x^3-7 x^2+2\left(m^2+6 m\right) x-8=0$.
	\shortans{-6} 
	 \loigiai{
		\begin{itemize}
			\item Điều kiện cần: Giả sử phương trình đã cho có ba nghiệm phân biệt $x_1,\, x_2,\, x_3$ lập thành một cấp  số nhân.\\
			Ta có 
			$x^3-7  x^2+2\left(m^2+6 m\right) x-8=\left(x-x_1\right)\left(x-x_2\right)\left(x-x_3\right), \forall m \in \mathbb{R} \Rightarrow x_1 x_2 x_3=8
			 $.
			Theo tính chất của cấp  số nhân $x_1 \cdot x_3=x_2^2$. \\
			Suy ra  $x_2^3=8 \Rightarrow x_2=2$.\\
			Thay  nghiệm $x=x_2=2$ vào phương trình đã cho, ta có
			$8-28+2\left(m^2+6 m\right) \cdot  2-8=0
			\Leftrightarrow  4 m^2+24 m-28=0
			 \Leftrightarrow\hoac{
			&m=1 \\ 
			&m=-7}$ 
			\item Điều kiện đủ: Thử lại  với các giá trị ${m}$ tìm được.\\
			Với $m=1$, ta có phương trình $x^3-7 x^2+14x-8=0 \Leftrightarrow \hoac{&x=4 \\ &x=2 \\  &x=1}$ (thoả mãn).\\
			Với $m=-7$, ta có phương  trình $x^3-7 x^2+14x-8=0 \Leftrightarrow \hoac{&x=4 \\ &x=2 \\ &x=1}$ (thoả mãn).\\
			Vậy $m=1;\, m=-7$ là  các giá trị cần tìm, suy ra $1+(-7)=-6$.
		\end{itemize}
	} 
\end{ex}


%%%%%%%%%%%%%%%% TỰ LUẬN %%%%%%%%%%%%%%%%

% nguồn: Ok_TuanLuan_Toan11_Lan2
\begin{ex}%[1D2V3-5]%[ma:70fe5446]
	Tìm $m$ để phương trình sau có  $3$ nghiệm lập thành cấp số nhân $x^3+(5-m) x^2+(6-5 m) x-6 m=0$.
	\loigiai{ 
		\[x^3+(5-m) x^2+(6-5 m)  x-6 m=0 \Leftrightarrow (x-m)(x^2+5x+6)=0\Leftrightarrow \hoac{&x=m\\&x=-2\\&x=-3.}\]
		Để 3 nghiệm lập thành  CSN xét 3 TH
		\begin{itemize}
			\item TH1: $-3 < m <-2 \Rightarrow m^2=6 \Leftrightarrow\left[\begin{array}{l}m=\sqrt{6} \\ m=-\sqrt{6}
		\end{array} \Rightarrow m=-\sqrt{6}\right.$.
		\item TH2: $-3  <-2 < m \Rightarrow 4=-3 m \Leftrightarrow m=-\dfrac{4}{3}$.
		\item TH3: $m <-3  <-2 \Rightarrow 9=-2 m \Leftrightarrow m=-\dfrac{9}{2}$.
	\end{itemize}
	Vậy có 3  giá trị của $m$ thỏa mãn.
} 
\end{ex}
